
Let $\Sigma=\{0,1\}$. In this section, we show how given a
quantifier-free formula $\Phi(x_1,\ldots,x_d)$ of Sem{\"e}nov
arithmetic, we can construct an \laVASS $V$ over the alphabet
$\Sigma_d\defeq \{0,1\}^d$ such that $\eval{L(V)}= \{\vec
x\in\N^d:\Phi(\vec x)\}$. We will subsequently observe that the
resulting \laVASS enjoy strong structural restrictions, giving rise to
the fragment of restricted \laVASS that we then formally define. For
our purposes, it will be sufficient to primarily focus on formulas
$\Phi$ of Sem{\"e}nov arithmetic which are conjunctions of atomic
formulas.
%
\begin{lemma}\label{lem:conjunctive-lavass}
  Consider a positive conjunctive formula of Sem{\"e}nov arithmetic
  \[
  \Phi(\vec x) \equiv A\cdot \vec x = \vec b \wedge
  \bigwedge_{i\in I} x_i = 2^{y_i},
  \]
  where $A\in \Z^{m\times n}$, $\vec b\in \Z^m$, $I$ is
  a finite index set, and $x_i$ and $y_i$
  are variables from $\vec x$. There is an \laVASS $V$ of dimension
  $2 \abs I$ and of size $(\norm A_{1,\infty} + \norm{\vec b}_\infty +
  2)^{O(m+\abs I)}$ such that $L(V)=\eval \Phi$.
\end{lemma}
%

We derive this lemma in two parts. First, it is well known that the
sets of solutions of the systems of linear equations $A\cdot \vec x
= \vec b$ can be represented by a regular language and are hence definable via
an \laVASS.
%\cite{WB00}, 
\begin{lemma}[\cite{WB00}, see also {\cite[Eqn.~(1)]{GuepinHW19}} ]\label{lem:equations-lavass}
  Given a system of equations $\Phi \equiv A\cdot \vec x = \vec b$
  with $A \in \Z^{m\times d}$ and $\vec b\in \Z^m$, there is a DFA $V$
  with at most $2^m \cdot \max\{ \norm A_{1,\infty}, \norm{\vec
  b}_\infty \}^m$ states such that $L(V)$ is zero-closed and
  $\eval{L(V)}=\eval \Phi$.
\end{lemma}
%

The crucial part, which requires the power of \laVASS, are exponential
equations $x=2^y$. An \laVASS $V$ with two counters and
$\eval{L(V)}=\eval{x=2^y}$ is depicted
in \Cref{fig:exponential-gadget}. Control-states are depicted as
circles and transitions as arrows between them. The vector before the
colon is the alphabet symbol read. For instance, the transition from
$q_0$ to $q_1$ reads the alphabet symbol $(1,0)\in \{0,1\}^2$. After a
colon, we display the counter operations when reading the alphabet
symbol, the operation on the first counter is displayed on the top and
the operation on the second counter on the bottom. Here and
subsequently, for presentational convenience, $\wID$ is the identity
$x\mapsto x$, and $\wTT$ and $\wTTPO$ are the functions $x \mapsto 2
x$ and $x \mapsto 2 x + 1$, respectively.  Thus, the transition from
$q_0$ to $q_1$ applies the identity function on the first counter, and
the function $x\mapsto 2x$ on the second counter.
%
\begin{figure}[!t]
\centering
\begin{tikzpicture}[->,shorten >=1pt,auto,node distance=4cm,semithick]

\tikzstyle{every state}=[draw=black,text=black,minimum size=25pt]

\node[initial by arrow,state,initial text=] (1) {$q_0$};
\node[state, accepting] (2) [right of=1] {$q_1$};

\path 
    (1) edge node [align=center]{\small{$\left[\begin{matrix} 1\\0 \end{matrix}\right] : \begin{array}{c}
         \wID  \\
         \wTT
    \end{array}$}}  (2)
    (1) edge [loop above] node {\small{$\left[\begin{matrix} 0\\0 \end{matrix}\right] : \begin{array}{c}
         \wID  \\
         \wTT
    \end{array}$}} (1)
    (2) edge [loop above] node {\small{$\left[\begin{matrix} 0\\b \end{matrix}\right] : \begin{array}{c}
         \wPP  \\
         \wTT + b
    \end{array}$}} (2);
\end{tikzpicture}
  \caption{Gadget with two counters for exponential equations $x=2^y$,
  where $b\in\{0,1\}$.}  \label{fig:exponential-gadget}
\end{figure}

%
The idea behind the gadget in \Cref{fig:exponential-gadget} is as
follows. For an example, suppose that $y=5$, then $x=32$, and in
binary the sequence of digits of $x$ and $y$ looks as follows:
%
\[
\left[
\begin{smallmatrix}
  x \\ y
\end{smallmatrix}
\right] = \left[
\begin{smallmatrix}
  1 \\ 0
\end{smallmatrix}
\right]
\left[
\begin{smallmatrix}
  0 \\ 0
\end{smallmatrix}
\right]
\left[
\begin{smallmatrix}
  0 \\ 0
\end{smallmatrix}
\right]
  \cdots
\left[
\begin{smallmatrix}
  0 \\ 0
\end{smallmatrix}
\right]
\left[
\begin{smallmatrix}
  0 \\ 1
\end{smallmatrix}
\right]
\left[
\begin{smallmatrix}
  0 \\ 0
\end{smallmatrix}
\right]
\left[
\begin{smallmatrix}
  0 \\ 1
\end{smallmatrix}
\right]
\]
%
Since $x=2^y$, we have that $x\in 0^*10^*$, and the number of tailing
zeros of $x$ is equal to the value of $y$. Thus, once a $1$ in the
binary representation of $x$ has been read, the first counter in the
gadget of \Cref{fig:exponential-gadget} keeps incrementing and counts
the number of tailing zeros of $x$. At the same time, the second
counter in the gadget of \Cref{fig:exponential-gadget} keeps the value
0 until it reads the first $1$ of the binary expansion of $y$, since
$2\cdot 0=0$. It then computes the value of $y$ in binary on the
second counter by multiplying the value of the second counter by 2
when reading a zero, and multiplying by two and adding one when
reading a one. The \laVASS in \Cref{fig:exponential-gadget} only
accepts when the first and the second counter have the same value,
i.e., when the number of tailing zeros of the binary expansion of $x$
equals the value of $y$, as required.
%
\begin{lemma}\label{lem:exponential-lavass}
  There is a fixed \laVASS $V$ of dimension two such that $L(V)$ is
  zero closed and $\eval{L(V)}=\eval{x = 2^y}$.
\end{lemma}
%

\Cref{lem:conjunctive-lavass} is now an easy consequence of
\Cref{lem:equations-lavass,lem:exponential-lavass} together with the
closure of \laVASS languages under intersection
(\Cref{prop:union-intersetion-closure}) and inverse homomorphisms
(\Cref{prop:projection-closure}).
